\chapter{Úvod}
\setcounter{page}{1}

Správa majetku je nedílnou součástí provozu každé organizace. Zahrnuje evidenci hmotných položek, sledování jejich umístění, stavu a~životního cyklu od pořízení přes inventarizaci až po vyřazení. V~mnoha organizacích, zejména ve školách a~menších institucích, je tato evidence vedena v~tabulkových procesorech bez systematické podpory inventarizace.

Tato práce vznikla na základě konkrétní potřeby Základní školy Ústavní v~Praze~8, která dosud evidovala majetek v~tabulkách aplikace Microsoft Excel. Při tisících položek rozmístěných v~desítkách místností je ruční inventura časově náročná a~náchylná k~chybám. Škola proto hledala řešení, které by umožnilo automatizovanou inventarizaci pomocí čtečky čárových kódů, generování etiket pro označení majetku a~přehlednou správu evidence prostřednictvím webové aplikace.

Cílem práce je navrhnout a~implementovat webový systém pro správu majetku s~mobilní aplikací pro offline inventuru v~terénu. Systém má umožnit evidenci majetkových karet, generování a~tisk etiket s~čárovými kódy, sběr inventurních dat pomocí průmyslové čtečky Zebra a~vyhodnocení inventury včetně generování reportů.

\medskip

Práce je strukturována následovně. Kapitola~\ref{chap:research} analyzuje existující přístupy ke~správě majetku a~zdůvodňuje vývoj vlastního řešení. Kapitola~\ref{chap:requirements} definuje funkční a~obecné požadavky a~popisuje případy užití. Kapitola~\ref{chap:technologies} srovnává technologie a~zdůvodňuje jejich volbu. Kapitola~\ref{chap:architecture} popisuje návrh architektury systému včetně datového modelu. Kapitola~\ref{chap:implementation} se zaměřuje na klíčové implementační detaily. Kapitola~\ref{chap:testing} pokrývá testování a~nasazení do produkčního prostředí. Kapitola~\ref{chap:conclusion} shrnuje dosažené výsledky a~navrhuje možnosti dalšího rozvoje.


